\chapter{Estado del arte}
\label{chap:estado-arte}

\section{De las anomalías a las variables explicativas}

Los retornos esperados de una acción no son independientes de sus características observables. Sobre
esa idea se construyó la literatura de factores: el tamaño (Banz, 1981), el valor (Fama y French,
1992, 1993), el momentum (Jegadeesh y Titman, 1993), la calidad (Asness et al., 2019), la
rentabilidad bruta (Novy-Marx, 2013) o la baja beta (Frazzini y Pedersen, 2014) quedaron asociados a
diferencias sistemáticas de rentabilidad. De ahí salen las cinco familias en que este trabajo
reparte sus variables ---calidad, valor, crecimiento, tendencia y riesgo---, y de ahí salen también
sus baselines: una fórmula determinista de un factor, calculada sobre el mismo panel, es el listón
mínimo que un sistema aprendido debe superar.

Esa literatura trae consigo una advertencia que importa más que el catálogo de primas. Las
regularidades documentadas se debilitan al medirse fuera de la muestra en que se descubrieron, y
más aún después de publicarse: McLean y Pontiff (2016) lo documentan sobre 97 predictores y Hou et
al. (2020) sobre 452 anomalías reevaluadas con criterios homogéneos, muchas de las cuales dejan de
superar umbrales exigentes o dependen de microcapitalizaciones. La implicación no es que toda
regularidad sea falsa, sino que una señal aprendida tiene que demostrar estabilidad temporal,
resistir controles conocidos y declarar cuánto de su éxito aparente procede de haber buscado entre
alternativas. Es el motivo de que este trabajo reserve una era completa y mida robustez antes de
afirmar nada.

Los factores cumplen aquí un tercer papel además de orientar el reparto y servir de baseline: son
controles de atribución. Si la señal desaparece al neutralizar exposiciones conocidas, la
interpretación debe ser más modesta. Conviene precisar, eso sí, que ningún agente representa un
factor puro: una empresa de alta calidad puede presentar a la vez baja volatilidad, crecimiento
persistente y valoración elevada. La especialización organiza la entrada del modelo, no impone que
el mercado recompense cada bloque por separado.

\section{Aprendizaje automático en \textit{asset pricing}}

Los métodos de aprendizaje automático amplían la capacidad de modelar no linealidades e
interacciones entre características. Gu et al. (2020) muestran que estos modelos pueden
mejorar la predicción transversal cuando el conjunto de señales es amplio y el entrenamiento respeta
la estructura temporal. LightGBM, usado aquí en cada agente, pertenece a la familia de árboles
potenciados: construye de forma secuencial árboles débiles que corrigen errores de los anteriores y
puede manejar relaciones no monotónicas (Ke et al., 2017).

Ese resultado forma parte de una literatura más amplia que intenta reducir dimensionalidad sin
suponer que cada característica actúa linealmente y de forma aislada. Kelly et al. (2019) vinculan
características y covarianzas mediante componentes principales instrumentados; Freyberger et al.
(2020) combinan selección regularizada y funciones no paramétricas; y Feng et al. (2020) estudian
qué factores aportan información incremental dentro de un catálogo amplio. Sus métodos son
distintos, pero convergen en dos ideas relevantes aquí: la no linealidad puede importar y añadir
predictores no garantiza añadir información.

La elección de árboles potenciados no se presenta por ello como una conclusión universal. Frente a
un modelo lineal, ganan capacidad para representar umbrales e interacciones, pero pierden una
interpretación directa de coeficientes. Frente a una red profunda, ofrecen una complejidad más
contenida y una atribución local más manejable para el tamaño del panel. El catálogo conserva una
alternativa lineal regularizada y una formulación de ranking directo precisamente para que estas
preferencias puedan contrastarse en vez de asumirse.

Hay además una diferencia entre predecir retornos y ordenar acciones. Un error pequeño en unidades
de retorno puede convivir con un orden transversal pobre, y una transformación monótona del score
puede alterar el error cuadrático sin cambiar ninguna decisión de selección. Este trabajo sigue la
segunda pregunta: cada agente aprende una representación continua, pero el protocolo juzga el orden
que induce. Esa elección acerca el objetivo estadístico al uso final sin hacer que la cartera entre
prematuramente en la selección.

Esa flexibilidad tiene un coste epistemológico. Un modelo flexible puede encontrar patrones espurios
con facilidad, especialmente si se prueban parámetros hasta maximizar una métrica final. Por ello la
pregunta no es si un algoritmo complejo puede ajustar la historia, sino si conserva una ventaja al
evaluarse en observaciones posteriores no utilizadas para su entrenamiento ni para su selección.
El diseño \textit{walk-forward}, las etiquetas cerradas y la era reservada son respuestas a esa
exigencia. La Tabla~\ref{tab:posicionamiento} enfrenta las cuatro líneas de literatura que este
trabajo cruza con el riesgo metodológico de cada una y con lo que aquí se hace al respecto.

\begin{table}[H]
\centering
\small
\caption{Posicionamiento metodológico del estudio frente a líneas de literatura relevantes.}
\label{tab:posicionamiento}
\begin{tabularx}{\textwidth}{lXXX}
\toprule
Enfoque & Aportación principal & Riesgo metodológico & Respuesta de este TFM\\
\midrule
Factores clásicos & Interpretabilidad económica y baselines simples & Confundir una prima histórica con señal universal & Comparación sobre el mismo panel y neutralización de estilo\\
ML transversal & Interacciones y no linealidades & Sobreajuste de parámetros y dependencia temporal & \textit{Walk-forward}, catálogo cerrado y Rank-IC OOS\\
Backtest de cartera & Traducción a utilidad económica & Elegir por CAGR, ignorar costes o supervivencia & Cartera posterior al ganador y costes explícitos\\
Inferencia con múltiples pruebas & Penalización del \textit{data snooping} & Publicar sólo la mejor especificación & Bootstrap, permutación, placebos y Deflated Sharpe\\
\bottomrule
\end{tabularx}
\caption*{Elaboración propia a partir de Gu et al. (2020), Harvey et al. (2016) y
López de Prado et al. (2014).}
\end{table}

\section{Sesgos que condicionan la evidencia}

El \textit{lookahead bias} aparece cuando una variable incorpora información que no era conocida en
la fecha de decisión. En fundamentales, usar la fecha de cierre fiscal en lugar de la fecha de filing
es un ejemplo clásico: el trimestre puede haber terminado, pero sus resultados no se publicaron
aún. El sesgo de supervivencia surge cuando se usa la composición actual de un índice para describir
el pasado; se eliminan así empresas que salieron del universo y se exagera la calidad de la muestra
(Brown et al., 1992).

El \textit{data snooping} es distinto: incluso con datos perfectamente cerrados, probar muchas
alternativas incrementa la probabilidad de encontrar una aparente ganadora. Harvey et al. (2016)
alertan del gran número de factores propuestos y de la necesidad de elevar el estándar de evidencia,
y de ahí salen las correcciones que este trabajo emplea ---el Deflated Sharpe Ratio de Bailey y
López de Prado (2014)--- y las que formalizan la misma preocupación por otras vías: el
\textit{Reality Check} de White (2000), el test de capacidad predictiva superior de Hansen (2005) o
la probabilidad de sobreajuste del backtest de Bailey et al. (2017). Ninguna reconstruye una
hipótesis verdaderamente previa después de observar los datos, y por eso el catálogo cerrado, el
registro de decisiones y la era reservada cumplen una función que ningún $p$-valor sustituye.

El sesgo de supervivencia tampoco se reduce a incluir o excluir empresas quebradas. Shumway (1997)
mostró que omitir retornos de exclusión puede distorsionar la medición. En este proyecto el problema
adopta otra forma: la pertenencia histórica se conoce, pero el proveedor gratuito deja de servir
precios de numerosos símbolos retirados. El universo correcto y el panel observable son, por tanto,
dos poblaciones diferentes. Esta distinción justifica que la cobertura se mida contra los miembros
reales del índice y no contra las filas que el propio pipeline consiguió construir.

\section{Medición de señal, amplitud y transferencia: Rank-IC como criterio de selección}

La métrica que selecciona el sistema es la correlación de rangos, no el error de una predicción
puntual. El motivo es práctico: para elegir acciones importa más acertar cuál de dos empresas irá
mejor que estimar con precisión cuánto rendirá cada una. Tiene además una propiedad cómoda: si un
modelo cambia de escala pero conserva el orden, la cifra no se mueve.

Una sola cifra, eso sí, no basta. Junto a ella se reporta en qué fracción de fechas fue positiva,
cuánto varía y su cociente entre media y desviación. Y hay una corrección obligada: como las
etiquetas a doce meses se solapan entre fechas consecutivas, el número de observaciones aparenta ser
mayor de lo que realmente es, lo que obliga a trabajar por bloques y con errores estándar
corregidos.

La literatura de gestión activa distingue además la calidad de una señal de la capacidad de
aprovecharla, y lo resume en una relación clásica ---la ley fundamental de Grinold (1989)--- que
este trabajo emplea como marco conceptual y no como identidad exacta. Dice que el resultado que un
gestor puede alcanzar depende de tres cosas: cómo de buena es su señal, cuántas decisiones
independientes puede tomar con ella, y qué parte de esa señal consigue llevar realmente a la
cartera.

\[
  \text{IR} \;\approx\; \text{IC} \times \sqrt{\text{amplitud}} \times \text{TC}
\]

Clarke et al. (2002) muestran por qué las restricciones reducen ese coeficiente de transferencia:
una señal puede contener información que una cartera \textit{long-only}, concentrada o sometida a
límites de rotación no consigue expresar. El punto es central en este TFM. El Rank-IC evalúa toda la
sección transversal, incluida la mitad inferior que una cartera sin posiciones cortas nunca puede
monetizar, mientras que la cartera mantiene ocho nombres y puede retenerlos doce meses. El salto
entre ambas métricas no es una conversión automática, sino un segundo problema de decisión.

Los costes añaden otra ruptura. Amihud (2002) relaciona iliquidez y retornos, y Novy-Marx y Velikov
(2016) muestran que el turnover puede consumir buena parte de las anomalías más intensivas en
negociación. Restar una comisión fija al final no basta cuando el coste también debe cambiar la
decisión de operar. Por eso esta memoria distingue la ruta congelada ---qué habría costado ejecutar
las mismas órdenes--- de la ruta resimulada, en la que un coste mayor eleva los umbrales y puede
evitar la operación.

La utilidad de esta descomposición es que separa tres causas de un mal resultado que una cifra de
rentabilidad confundiría: que la señal no ordene, que se tomen pocas decisiones independientes o que
la cartera no consiga expresar lo que la señal indica. Las tres restricciones de este diseño actúan
sobre el tercer término: la cartera es \textit{long-only}, de modo que la mitad de la información
que el Rank-IC mide es por construcción inaccesible; mantiene unas pocas posiciones sobre un
universo de unos quinientos valores; y paga costes proporcionales a su rotación. De ahí que la
métrica de selección sea el Rank-IC y no el CAGR: seleccionar por rentabilidad final confundiría los
tres términos, mientras que medir la señal aislada permite atribuir después cuánto pierde la
implementación.

Hay una quinta línea que no es metodológica sino estructural, y que explica buena parte de lo
anterior: la aritmética de la gestión activa de Sharpe (1991), según la cual el conjunto de los
gestores activos obtiene, antes de costes, exactamente el rendimiento del índice. Un trabajo que
reporta alfa persistente sobre un índice amplio está afirmando, implícitamente, que ese alfa lo pagó
otro participante. Es la razón de fondo por la que la replicación encuentra tanta caída fuera de
muestra: una anomalía publicada deja de ser una ventaja informativa en cuanto el resto del mercado
puede explotarla. La distinción entre calidad de señal e implementación que formula López de Prado
(2018) es la que este trabajo adopta como orden de trabajo: primero la señal, después la cartera, y
nunca al revés.
